\section{Related Work}

Our work intersects three research threads: output-level uncertainty quantification, hidden state probing, and transformer interpretability. We position each as necessary but insufficient for efficient correctness prediction.

\subsection{Output-Level Uncertainty Quantification}

\textbf{Token-based methods.} \citet{fadeeva2024survey} survey token-level uncertainty metrics including entropy, perplexity, and sequence probability. These methods achieve $\sim$0.62--0.65 AUROC for correctness prediction on factual QA---better than random but insufficient for reliable deployment. The fundamental limitation: output softmax distributions conflate linguistic confidence (fluency) with factual confidence (correctness).

\textbf{Semantic entropy.} \citet{kuhn2023semantic} address this by sampling multiple responses and clustering by semantic equivalence. Semantic entropy achieves $\sim$0.80 AUROC by detecting when diverse samples produce semantically distinct answers (high uncertainty) versus consistent answers (low uncertainty). However, this requires 5--20 forward passes per question, making it impractical for real-time applications.

\textbf{Single-sample approximations.} Recent work attempts to approximate semantic entropy from single samples. \citet{kossen2024semantic} train probes to predict semantic entropy from hidden states, achieving 0.85+ correlation with ground-truth SE. However, predicting \emph{uncertainty} differs from predicting \emph{correctness}---a confident model can still be wrong.

\subsection{Hidden State Probing}

\textbf{Probing for concepts.} Linear probes have successfully extracted semantic concepts from transformer hidden states, including sentiment, syntax, and factual knowledge. \citet{burns2023discovering} use contrastive probing to detect model beliefs. Our work extends this paradigm: rather than probing for specific facts, we probe for correctness of model-generated answers.

\textbf{Semantic Entropy Probes (SEP).} \citet{kossen2024semantic} is closest to our approach. SEP trains linear probes on hidden states to predict semantic entropy, enabling single-pass uncertainty estimation. We depart from SEP in two ways: (1) our target is ground-truth correctness, not uncertainty; (2) we systematically characterize layer depth effects, identifying optimal extraction at 50--60\% depth.

\subsection{Transformer Interpretability}

\textbf{Logit lens.} \citet{nostalgebraist2020logit} demonstrates that projecting intermediate hidden states to vocabulary space reveals semantic content developing across layers. This supports our hypothesis that middle layers encode semantic knowledge before output formatting.

\textbf{Layer-wise analysis.} \citet{aiersilan2026hallucination} study hallucination detection across layers, finding optimal signals at 40--56\% depth for Llama and Mistral models. Our inverted-U pattern (peak at 50\% depth) aligns with these findings, confirming middle layers as the locus of semantic representations.
